\section{Problem Setting and Threat Model}
\label{sec:threat}

\subsection{Pipeline and derived artifacts}
A request from user $u$ passes through a semantic-cache lookup, retrieval of the top-$k$ chunks from
a vector index, reranking, and generation. The application keeps per-session chat memory: recent
turns verbatim and, past a length threshold, a rolling summary written by the generator. The
identity system (IAM) is the source of truth. Three artifacts hold derived state that is not
recomputed per request:
\begin{itemize}
\item the \emph{index ACLs}, a copy of document permissions stored beside the embeddings and
refreshed by a periodic crawl, by webhook, or both;
\item the \emph{semantic cache}, which returns a stored answer when a new question's embedding is
within a similarity threshold, until a time-to-live (TTL) expires;
\item the \emph{chat memory}, whose summary condenses earlier answers and names no source document.
\end{itemize}
\Cref{tab:grounding} grounds each parameter in documented framework and product behaviour.

\begin{table*}[t]
\centering
\caption{Grounding of the deployment model in documented behaviour (documentation accessed October
2026). The \texttt{baseline} preset uses the documented defaults; settings that deviate are named
presets of their own.}
\label{tab:grounding}
\footnotesize
\setlength{\tabcolsep}{3pt}
\begin{tabular}{@{}p{2.6cm}p{6.6cm}p{7.4cm}@{}}
\toprule
Component & Documented behaviour & Our setting \\
\midrule
Index ACLs & Security trimming stores user or group identifiers in an index field and filters on
them at query time~\cite{doc_azure_trimming}; vector stores filter on payload or metadata
fields~\cite{doc_qdrant,doc_pinecone,doc_chroma}; connectors ingest ACLs with items and re-crawl on a
schedule~\cite{doc_graph_connectors}. & Per-user ACL booleans in Chroma metadata, pre-filter in the
query; sync by hourly or daily crawl, webhook, or webhook with loss and a daily backstop. \\
Semantic cache & Similarity threshold and TTL; scoping by user or other fields is configured by the
application~\cite{doc_gptcache,doc_redisvl,bang2023gptcache,gill2024meancache}. & Threshold 0.90,
TTL 3600\,s; per-user scope in \texttt{baseline}, tenant-shared scope as the named misconfiguration
\texttt{baseline\_shared}. \\
Chat memory & Recent turns kept verbatim, older turns replaced by a model-written running
summary~\cite{doc_langchain_summary,doc_llamaindex_memory,wang2023recursively}. & Summary after 8
turns, 4 kept verbatim, written by the generator model; session retention swept (none, 1\,h, 24\,h). \\
Revocation & Access must be modified or disabled when users change role or
leave~\cite{nist80053r5}. & Single document, role removal, department move (mover), offboarding,
document reclassification. \\
\bottomrule
\end{tabular}
\end{table*}

\subsection{Adversaries and goals}
\emph{Revocation leakage.} A legitimate user has used a document and then loses access to it. The
adversary is that user acting as an ordinary authenticated client, or a colleague in the same
tenant who never had access. Access is removed for one document (single-document), by removing a
group (role), by a department transfer that removes one group and adds another (mover), or for the
whole account (offboarding), or when an administrator narrows the document's ACL
(reclassification), which removes it from every user of a group at once. A mover keeps an active session and account, which makes it the
realistic case; offboarding assumes the deprovisioned user can still reach the application (for
example until an SSO session or an API token expires) and is reported for completeness. The goal
is to learn secrets of the revoked document that the user was never shown (\emph{new disclosure}),
including those of a version written after the revocation, or to be shown again what they were
shown before (\emph{re-disclosure}). Re-disclosure gives the user nothing new, but it violates
complete mediation~\cite{saltzer1975protection} and the obligation to cut off access on a role
change~\cite{nist80053r5}, and it is the step that moves revoked content into state other users
read.

\emph{Existence inference.} An API client without access to a target document decides whether it
exists in the collection. An \emph{insider} belongs to the target's tenant; an \emph{outsider} does
not. The adversary sees what an API client sees: the response text, the latency and, depending on
the exposure profile, cited sources with relevance scores and the reranker confidence. It knows the
target's topic and a question about it (the membership-inference setting~\cite{shokri2017membership});
discovering unknown topics is out of scope.

\subsection{Threats}
\Cref{tab:threats} lists the leak mechanisms and the switch that addresses each.

\begin{table}[t]
\centering
\caption{Threats and the mitigation switch addressing each.}
\label{tab:threats}
\footnotesize
\begin{tabular}{@{}p{0.5cm}p{3.6cm}p{3.6cm}@{}}
\toprule
 & Mechanism & Mitigation \\
\midrule
T1 & Index ACL copy is stale & live IAM check \\
T2 & Cached answer outlives permission, or is shared across users & per-user scope; eviction, lazy taint or epoch check \\
T3 & Memory and summaries retain content & taint purge and lazy check; epoch check; flush \\
T4 & Retrieval score of a withheld candidate & ACL pre-filter; scores over authorized chunks; score bands \\
T5 & Reranker confidence & as T4 \\
T6 & Latency differs between paths & padding, deadline, no fast path \\
T7 & Refusal wording differs & uniform refusal text \\
T8 & A revoked user re-enters content received earlier, which enters shared state & query fingerprinting \\
\bottomrule
\end{tabular}
\end{table}

T1--T3 are revocation threats, evaluated by the revocation experiment. T4--T7 are existence channels,
evaluated by the twin-world experiment; T2 also carries existence signal when a cache is shared
between users. T8 does not arise from a copy the application made but from the user re-supplying
content, which is why assumption~A1 (\cref{sec:soundness}) states it and \cref{sec:res-launder}
measures it.
